\begin{helper}
Let \(\beta=1/2\), and let
\[
m(n)=M(n)=\noderef{TwoBiteNaturalReducedVertexCount}(n).
\]
Put \(p(n)=\noderef{TwoBiteEdgeProbability}(\beta,n)\).  Under
\(\noderef{TwoBiteEventProbability}(n,m(n),p(n))\), with high probability in
the sense of \(\noderef{WithHighProbability}\), every distinct red pair and
every distinct blue pair has lifted-neighborhood intersection size at most
\[
100(\log n)^3.
\]
Equivalently, with high probability,
\[
\begin{aligned}
&\forall r\ne s\in \operatorname{Fin}(m(n)),\quad
\left|
\noderef{TwoBiteLiftedNeighborhood}(\omega,\operatorname{inl}(r))
\cap
\noderef{TwoBiteLiftedNeighborhood}(\omega,\operatorname{inl}(s))
\right|
\le 100(\log n)^3,\\
&\forall b\ne c\in \operatorname{Fin}(m(n)),\quad
\left|
\noderef{TwoBiteLiftedNeighborhood}(\omega,\operatorname{inr}(b))
\cap
\noderef{TwoBiteLiftedNeighborhood}(\omega,\operatorname{inr}(c))
\right|
\le 100(\log n)^3 .
\end{aligned}
\]
\end{helper}
\begin{proof}
Write \(L=\log n\), \(M=M(n)\), and
\[
p=p(n)=\frac12\sqrt{\frac{L}{n}}.
\]
By \(\noderef{TwoBiteNaturalReducedVertexCount}\), together with
\(\noderef{TwoBiteReducedVertexCount}\) and \(\noderef{TwoBiteStretch}\), for
all sufficiently large \(n\),
\[
M=\left\lfloor\frac n{L^2}\right\rfloor,\qquad
\frac12\frac n{L^2}\le M\le\frac n{L^2},\qquad n\le M^2.
\tag{1}
\]
The support and normalization facts needed for the injection coordinate at
this rounded size are supplied eventually by
\(\noderef{TwoBiteNaturalTotalMassOneEventually}\), and
\(\noderef{OppositeProjectionEdgeProbBounds}\) gives \(0\le p\le1\) after
discarding finitely many \(n\).  These finite initial values do not affect
the high-probability limit.  From \(M\le n/L^2\) and the displayed formula
for \(p\),
\[
Mp^2\le \frac1{4L},\qquad M^2\le n^2=\exp(2L).
\tag{2}
\]

We first record the sharper same-color base common-neighbor estimate needed
for this lifted-intersection slice.  Fix distinct red vertices
\(r,s\in\operatorname{Fin}(M)\).  By \(\noderef{TwoBiteSample}\), an outcome
is a triple consisting of a red graph \(G_R=\noderef{TwoBiteRedGraph}(\omega)\),
a blue graph \(G_B=\noderef{TwoBiteBlueGraph}(\omega)\), and an injection of
the \(n\) final vertices into the \(M^2\) product cells.  If \(A\) is any
event depending only on \(G_R\), then
\(\noderef{TwoBiteEventProbability}\) writes its probability as the finite
sum of \(\noderef{TwoBiteSampleWeight}\) over all triples satisfying \(A\).
The summands are nonnegative by \(\noderef{TwoBiteSampleWeightNonnegative}\)
and \(0\le p\le1\).  The definition of
\(\noderef{TwoBiteSampleWeight}\) factors each summand as
\[
\noderef{GnpGraphWeight}(p,G_R)\,
\noderef{GnpGraphWeight}(p,G_B)\,
\noderef{UniformInjectionWeight}(n,M).
\tag{3}
\]
Because \(A\) ignores \(G_B\) and the injection, the sum over blue graphs is
\(1\) by \(\noderef{GnpGraphWeightSumOne}\), and the sum over injections is
\(1\) because \(n\le M^2\) and
\(\noderef{UniformInjectionWeight}(n,M)\) is uniform on the nonempty set of
injections.  Hence the red marginal under the two-bite law is exactly the
\(G(M,p)\) product law described by \(\noderef{GnpGraphWeight}\).  Reversing
red and blue gives the same product-law marginal for
\(\noderef{TwoBiteBlueGraph}\).

For \(t\notin\{r,s\}\), the edge indicators of \(\{r,t\}\) and \(\{s,t\}\)
are independent Bernoulli \(p\) variables under this product law, and for
different \(t\)'s all involved edge coordinates are distinct.  Thus the red
base common-neighbor count
\[
a_{r,s}=\left|\{t:\,G_R.Adj\,r\,t\ \text{and}\ G_R.Adj\,s\,t\}\right|
\]
has distribution \(\operatorname{Bin}(M-2,p^2)\), with mean at most
\(Mp^2\le 1/(4L)\).  The identical statement holds for each distinct blue
pair \(b,c\).

Use the elementary counting tail bound
\[
\Pr(Y\ge T)\le \left(\frac{e\mu}{T}\right)^T
\tag{4}
\]
for a binomial or hypergeometric count \(Y\) with mean \(\mu\): if \(Y\ge T\),
some \(T\)-set of candidate coordinates is entirely marked, and the union
bound over \(T\)-sets gives
\(\binom{N}{T}q^T\le(eNq/T)^T=(e\mu/T)^T\).  Apply (4) with
\(T=\lceil L\rceil\) to the variable \(a_{r,s}\).  Since
\(\mu\le1/(4L)\), for all large \(n\),
\[
\Pr(a_{r,s}>L)\le
\left(\frac{e/(4L)}{L}\right)^L
\le \exp(-{\tfrac12}L\log L).
\tag{5}
\]
There are fewer than \(M^2\le\exp(2L)\) ordered red pairs and the same number
of ordered blue pairs.  The union bound and (5) imply that, with probability
tending to \(1\), the event \(\mathcal B_n\) holds: all distinct same-color
base pairs, red and blue, have common-neighbor count at most \(L\).

Now condition on the red base graph and assume \(\mathcal B_n\).  Fix a
distinct red pair \(r,s\), and set
\[
a=a_{r,s}\le L.
\]
By the definitions of \(\noderef{TwoBiteLiftedNeighborhood}\),
\(\noderef{TwoBiteRedProjection}\), and \(\noderef{TwoBiteEmbedding}\), a
final vertex belongs to both lifted neighborhoods
\[
\noderef{TwoBiteLiftedNeighborhood}(\omega,\operatorname{inl}(r))
\quad\text{and}\quad
\noderef{TwoBiteLiftedNeighborhood}(\omega,\operatorname{inl}(s))
\]
exactly when its injected product cell has red coordinate in
\(N_{G_R}(r)\cap N_{G_R}(s)\).  Therefore the marked product cells are
\[
(N_{G_R}(r)\cap N_{G_R}(s))\times\operatorname{Fin}(M),
\]
and this marked set has cardinality \(aM\).

Conditional on the fixed red graph, the blue graph weights again normalize
to \(1\) by \(\noderef{GnpGraphWeightSumOne}\), and
\(\noderef{UniformInjectionWeight}\) makes the injection a uniform ordered
sample of \(n\) distinct cells from the \(M^2\) product cells.  Hence the
lifted intersection size
\[
H_{r,s}=
\left|
\noderef{TwoBiteLiftedNeighborhood}(\omega,\operatorname{inl}(r))
\cap
\noderef{TwoBiteLiftedNeighborhood}(\omega,\operatorname{inl}(s))
\right|
\]
is hypergeometric with population size \(M^2\), marked-set size \(aM\), and
sample size \(n\).  Its conditional mean is
\[
\mu_{r,s}=n\frac{aM}{M^2}=\frac{na}{M}
\le \frac{nL}{M}\le 2L^3,
\tag{6}
\]
where the last inequality uses (1).  This is the same without-replacement
model controlled by \(\noderef{HypergeometricMgfComparison}\), and the
counting bound (4) also applies directly to it because \(H_{r,s}\ge T\)
requires some \(T\) of the sampled coordinates to lie in the marked set.
With
\[
T=\lceil 100L^3\rceil,
\]
the event \(H_{r,s}>100L^3\) implies \(H_{r,s}\ge T\), and (4), (6) give
\[
\Pr(H_{r,s}>100L^3\mid G_R,\mathcal B_n)
\le \left(\frac{2eL^3}{100L^3}\right)^{100L^3}
\le \exp(-c_1L^3)
\tag{7}
\]
for a fixed constant \(c_1>0\) and all sufficiently large \(n\).

The blue case is identical.  Conditioning on
\(G_B=\noderef{TwoBiteBlueGraph}(\omega)\), the marked cells for a distinct
blue pair \(b,c\) are
\[
\operatorname{Fin}(M)\times (N_{G_B}(b)\cap N_{G_B}(c)),
\]
and \(\noderef{TwoBiteBlueProjection}\) identifies the blue lifted
intersection with the final vertices injected into this marked set.  On
\(\mathcal B_n\), its marked-cell count is at most \(LM\), so the same
hypergeometric bound (7) holds.

There are fewer than \(M^2\) ordered red pairs and fewer than \(M^2\) ordered
blue pairs.  On \(\mathcal B_n\), the union bound gives conditional failure
probability at most
\[
2M^2\exp(-c_1L^3)
\le 2\exp(2L-c_1L^3)\to0.
\tag{8}
\]
Since \(\Pr(\mathcal B_n)\to1\), (8) proves that, with high probability, all
red-red and blue-blue lifted-neighborhood intersections are bounded by
\(100(\log n)^3\).  The mixed red-blue intersection estimate is a different
slice and is not included here.
\end{proof}
